\section{征发失去回音}
\label{sec:sui-collapse}

大业七年王薄在长白山起事，表明山东的逃亡和不服从已能取得持续武装的形态；次年征辽失败则把兵役、粮运与地方治安的压力进一步推向河北、山东。把民变直接归结为一首民歌、或归结为炀帝的某一种性格，都会掩盖中间的行政过程：人要被登记、征发、输送，也要在道路、山泽和地方网络中决定是否逃离。现存材料能确认起事与镇压，却很难逐户说明参加者的动机。%
\footnote{王薄起事见 ST-611-WANG-BO-REBELLION，征辽失利见 ST-612-GOGURYEO-INVASION，依据《隋书·炀帝纪下》《隋书·王薄传》《隋书·东夷传·高丽》与《资治通鉴》隋纪。民歌、兵额及参与规模都经后世文本保存，不能视为统计。}

613年杨玄感利用皇帝在辽东、自己又掌黎阳转运的空隙起兵；615年始毕可汗围雁门，炀帝须靠勤王与外交解围。前者说明仓储、转运与重臣任用可以转化为政变资源，后者说明征辽抽调的军力无法同时保证北边的安全。两事都不是孤立的宫廷事故：当中央必须不断请求地方兵力来填补缺口时，地方军队和豪强也获得了更高的政治位置。%
\footnote{杨玄感起兵见 ST-613-YANG-XUANGAN，雁门之围见 ST-615-YANMEN-SIEGE，依据《隋书·炀帝纪下》《隋书·杨玄感传》《隋书·突厥传》《北史》与《资治通鉴》。同谋范围、外交传讯与围城细节存在叙事差异。}

大业十四年江都兵变中，宇文化及等杀炀帝并挟隋廷北返，表明朝廷甚至已不能稳定供养和指挥自己的骁果军。皇帝遇害不是“天下”立即转换的按钮，却使各地竞争者更容易以不同国号争夺名义。隋的崩解因而应理解为征发、战争、供给、边防和军队忠诚一起断裂的过程；唐初史书所提供的亡国图式可以帮助定位事件，却不能替代这些多重条件。%
\footnote{江都兵变见 ST-618-JIANGDU-COUP，依据《隋书·炀帝纪下》《隋书·宇文化及传》及《资治通鉴》隋纪。参与者责任、遗诏和宫廷情节多带唐初道德化叙事。}
